\section{The ordinary-sofa bound used in the proof}\label{app:ordinary-sofa}

An ordinary moving sofa is a nonempty compact connected planar set
with a passage through $\Lm$ as in \cref{def:passage}. Baek's
optimality theorem for Gerver's sofa supplies the following bound.

\begin{theorem}[Baek \cite{Baek}]\label{thm:ordinary-sofa-bound}
Every ordinary moving sofa $A$ satisfies
\[
 |A|\le |S_{\rm G}|\le2.2199,
\]
where $S_{\rm G}$ is Gerver's sofa.
\end{theorem}

We use this theorem as an area obstruction. Each application first
constructs a connected survivor with an explicit ordinary passage, then
shows that its area would exceed the bound. The exact value of
$|S_{\rm G}|$ is immaterial here. We use the
rational relaxation $G_0=2.2199$ in two strict estimates for
full-turn caps. No classification of ordinary equality cases enters
the argument.

\subsection{Why a cap survivor is an ordinary sofa}

For a height-one downward cap $U$, let
\[
 A_U=\{(x,y)\in U:n_{\pi/2}(U)(x)\le y\}.
\]
An arbitrary cap need not leave a connected survivor: its niche can
reach above its roof. The proof therefore establishes the following
geometric condition before applying the ordinary bound.

\begin{lemma}\label{lem:ordinary-survivor}
Suppose $U$ is a downward cap of height one and
\[
 n_{\pi/2}(U)(x)\le f_U(x)\quad(x\in I_U).
\]
Then $A_U$ is a nonempty compact connected ordinary moving sofa.
Its area is
\[
 |A_U|=|U|-\int_{I_U}n_{\pi/2}(U)(x)\,\dd x.
\]
\end{lemma}
\begin{proof}
The niche is continuous by \cref{lem:barrier-continuity}, so the
survivor is closed in $U$ and compact. Each of its fibers is the
nonempty interval $[n_{\pi/2}(U)(x),f_U(x)]$. The concave roof
is continuous on the interior of $I_U$, and is continuous at its
endpoints as well: concavity gives the lower limit there, while
compactness gives the upper limit. Its graph is consequently
connected. Every point of $A_U$ joins this graph by a vertical
segment in $A_U$, proving connectedness. The area formula follows
by integrating these fiber lengths.

Translate $U$ horizontally so that its right endpoint is $1$.
Write $F(t)=h_U(t)$ and $G(t)=h_U(t+\pi/2)$, and use
\[
 p\longmapsto R_{-t}p+(1-F(t),1-G(t)),\qquad
                 0\le t\le\pi/2.
\]
The outer walls contain $U$ by definition of the support function.
At every interior angle, the inequality $y\ge n_{\pi/2}(U)(x)$
places every survivor point on or above at least one of the two
inner walls. Thus the placement lies in $\Lm$. At $t=0$,
$F(0)=G(0)=1$, so the initial placement is the identity and lies
in $\Hin$. At $t=\pi/2$, its first coordinate is $y$, in $[0,1]$,
and its second is $1-h_U(\pi)-x\le1$. Hence it lies in $\Vm$.
The support function is continuous, so this is a continuous passage.
\end{proof}

The full-turn estimates verify the displayed roof inequality for
the caps at which $G_0$ is used. They retain the complete niche,
including any part outside the middle half. This distinguishes the
area in \cref{lem:ordinary-survivor} from the spatial score
$P_{\pi/2}$, which charges only the exterior roof and the middle
barrier. The proof converts between them by explicit exterior and
niche terms before applying \cref{thm:ordinary-sofa-bound}.
